\section*{Chapter \ref{ch:b1:diagenesis-basins} --- Diagenesis and Sedimentary Basins}

\begin{solution}{exo:b1:diagenesis-basins:1}
Sandstone $0.49\,e^{-2000/3704} = 0.49\times0.583 = \num{0.286}$; shale
$0.63\,e^{-2000/1960} = 0.63\times0.360 = \num{0.227}$. The porosity halves when
$e^{-z/\lambda} = 1/2$, at $z = \lambda\ln 2$: \qty{2567}{m} for sandstone,
\qty{1359}{m} for shale.
\end{solution}

\begin{solution}{exo:b1:diagenesis-basins:2}
$\Delta A = \Delta E + S = 30 + 50 = \qty{80}{m}$. Of it, \qty{60}{m} is filled
with sediment, so $\Delta W = 80 - 60 = \qty{20}{m}$: the shelf deepens by
\qty{20}{m}.
\end{solution}

\begin{solution}{exo:b1:diagenesis-basins:3}
$h = Y(\rho_m - \rho_w)/(\rho_m - \rho_s)$: with \qty{2300}{kg.m^{-3}},
$1.0\times2270/1000 = \qty{2.27}{km}$; with \qty{2500}{kg.m^{-3}},
$1.0\times2270/800 = \qty{2.84}{km}$. Each metre of sediment takes the place of
a metre of \omterm{def:b1:layered-earth:layers}{mantle} in the column; the denser the sediment, the smaller the mass
it lacks compared with the \omterm{def:b1:layered-earth:layers}{mantle}, so the more of it is needed to make up the
same deficit as the water-filled hole.
\end{solution}

\begin{solution}{exo:b1:diagenesis-basins:4}
(a)~Calcite $100.09/2.71 = \qty{36.93}{cm^3.mol^{-1}}$; dolomite $184.40/2.84
= \qty{64.93}{cm^3.mol^{-1}}$. (b)~Two moles of calcite occupy
\qty{73.86}{cm^3}, one of dolomite \qty{64.93}{cm^3}: the solid shrinks by
$1 - 64.93/73.86 = \qty{12.1}{\%}$. (c)~In a cubic metre, the \qty{0.90}{m^3}
of calcite becomes $0.90\times(1 - 0.121) = \qty{0.791}{m^3}$ of dolomite; the
porosity becomes $1 - 0.791 = \num{0.21}$.
\end{solution}

\begin{solution}{exo:b1:diagenesis-basins:5}
(a)~$s = 1000 - 0.63\times1960\times(1 - e^{-1000/1960}) = \qty{506.5}{m}$; the
mean porosity is $1 - 506.5/1000 = \num{0.493}$. (b)~With the top at
\qty{2000}{m}, solving $s(2000, b) = 506.5$ gives $b = \qty{2628.7}{m}$: the
layer is \qty{628.7}{m} thick, of mean porosity $1 - 506.5/628.7 = \num{0.194}$.
(c)~The grains are the same, so the water lost is the thickness lost:
$1000 - 628.7 = \qty{371}{m^3}$ of water per square metre.
\end{solution}

\begin{solution}{exo:b1:diagenesis-basins:6}
(a)~$s = 300 - 0.49\times3704\times(e^{-2500/3704} - e^{-2800/3704}) =
\qty{228.10}{m}$. (b)~From $b_0 = 300$: $s(0, 300) = 158.80$, $1 - \varphi(300) =
0.5481$, $b_1 = 300 + (228.10 - 158.80)/0.5481 = \qty{426.44}{m}$; then
$s(0, 426.44) = 229.07$, $1 - \varphi = 0.5633$, $b_2 = 426.44 - 0.96/0.5633 =
\qty{424.73}{m}$; a third step changes it by less than a millimetre. The
sandstone was \qty{424.7}{m} thick at the surface. (c)~$424.7/300 = 1.42$: burial
has thinned it by a factor of \num{1.42}.
\end{solution}

\begin{solution}{exo:b1:diagenesis-basins:7}
(a)~The ice holds $52\times10^6\times916.7/1000 = \qty{4.77e7}{km^3}$ of water;
spread over \qty{361.84e6}{km^2} it is $\qty{0.1317}{km} = \qty{132}{m}$ of sea
level. (b)~The shorelines give \qty{134}{m}, \qty{2}{m} more. (c)~$4.767\times
10^7/0.134 = \qty{355.7e6}{km^2}$, \qty{1.7}{\%} less than today. As sea level
fell the continental shelves emerged, so the glacial ocean was smaller, and the
same volume of water taken from it lowered it more.
\end{solution}

\begin{solution}{exo:b1:diagenesis-basins:8}
(a)~$(-9.1 - (-98.1))/8 = \qty{11.1}{m}$ per thousand years, close to the
\qty{12}{m} found from shorelines. (b)~\qty{40}{mm} a year is \qty{40}{m} per
thousand years, $40/12 = 3.3$ times the average; in five hundred years,
$0.040\times500 = \qty{20}{m}$.
\end{solution}

\begin{solution}{exo:b1:diagenesis-basins:9}
(a)~$S = 50\times2 = \qty{100}{m}$ and $\Delta E = \qty{-40}{m}$:
$\Delta A = \qty{60}{m}$, all filled with sediment since the depth is unchanged.
(b)~In \num{20000} years the basement subsides only $50\times0.02 = \qty{1}{m}$:
the \omterm{def:b1:diagenesis-basins:eustasy}{relative sea level} falls by $120 - 1 = \qty{119}{m}$. The sea floor, which lay
\qty{20}{m} below sea level, ends $119 - 20 = \qty{99}{m}$ above it: the whole
shelf emerges.
\end{solution}

\begin{solution}{exo:b1:diagenesis-basins:10}
(a)~For a thin layer the porosity is nearly constant across it: its solid
thickness is $s \approx h(1 - \varphi(z))$ at depth and $s \approx h_0(1 -
\varphi_0)$ at the surface. Equating the two gives $h_0 \approx h(1 -
\varphi(z))/(1 - \varphi_0)$. (b)~$\varphi(3000) = 0.63\,e^{-3000/1960} =
0.136$, so $h_0 \approx 10\times0.864/0.37 = \qty{23.3}{m}$, against
\qty{23.1}{m} exactly. (c)~Restored to the surface, the layer is
\qty{23}{m} thick, no longer thin compared with the decay of porosity near the
surface: its porosity falls from $\varphi_0$ at its top to less at its base, so
its grains fit in a little less than the thickness computed with $\varphi_0$
throughout.
\end{solution}

\begin{solution}{exo:b1:diagenesis-basins:11}
(a)~$s = 3000 - 0.63\times1960\times(1 - e^{-3000/1960}) = \qty{2032}{m}$; the
mean porosity is $1 - 2032/3000 = \num{0.323}$ and the mean density
$0.323\times1030 + 0.677\times2700 = \qty{2161}{kg.m^{-3}}$. (b)~The factor is
$(3300 - 1030)/(3300 - 2161) = \num{1.99}$; the fill represents a water-loaded
\omterm{def:b1:diagenesis-basins:subsidence}{subsidence} of $3000/1.99 = \qty{1505}{m}$. (c)~With the grain density the factor
would be $2270/600 = \num{3.78}$ and the \omterm{def:b1:diagenesis-basins:subsidence}{subsidence} only \qty{793}{m}. It would be
wrong because a third of the fill is water: the sediment is far lighter than its
grains, and only the density of the saturated sediment counts in the balance of
masses.
\end{solution}

\begin{solution}{exo:b1:diagenesis-basins:12}
(a)~The cube holds \qty{0.90}{m^3} $= \qty{9.0e5}{cm^3}$ of calcite, that is
$9.0\times10^5/36.93 = \num{24370}$ moles. (b)~One magnesium per two calcite:
\num{12180} moles. (c)~A kilogram of seawater holds $35.165\times0.0365 =
\qty{1.284}{g}$ of magnesium, $1.284/24.305 = \qty{0.0528}{mol}$; the reaction
needs $12180/0.0528 = \qty{2.31e5}{kg}$ of seawater, $\qty{224}{m^3}$ --- more
than two hundred times the volume of the rock, even if every ion of magnesium is
taken. (d)~Dolomite can form on a large scale only where very large volumes of
water flow through the rock for a long time, or where the water is a brine
richer in magnesium than seawater: in platforms flushed by sinking brines, for
instance, rather than in a sediment that keeps the same pore water.
\end{solution}

\begin{solution}{pb:b1:diagenesis-basins:1}
\textbf{1.} Sandstone: $\varphi(0) = \num{0.490}$, $\varphi(800) =
0.49\,e^{-800/3704} = \num{0.395}$. Shale: $\varphi(800) = \num{0.419}$,
$\varphi(2600) = \num{0.167}$. Chalk: $\varphi(2600) = \num{0.110}$,
$\varphi(3200) = \num{0.072}$.

\textbf{2.} With $s$ from \cref{prop:b1:diagenesis-basins:solid} and $\bar\varphi
= 1 - s/h$ (question~7): sandstone \num{0.441}, shale \num{0.274}, chalk
\num{0.090}.

\textbf{3.} $h\bar\varphi$: sandstone \qty{353}{m^3}, shale \qty{493}{m^3}, chalk
\qty{54}{m^3} of water per square metre.

\textbf{4.} $\bar\varphi\rho_w + (1 - \bar\varphi)\rho_g$: sandstone
$0.441\times1030 + 0.559\times2650 = \qty{1936}{kg.m^{-3}}$, shale
\qty{2242}{kg.m^{-3}}, chalk \qty{2559}{kg.m^{-3}}. The column weighs
$1936\times800 + 2242\times1800 + 2559\times600 = \qty{7.12e6}{kg}$ per square
metre: a mean density of \qty{2225}{kg.m^{-3}}.

\textbf{5.} Lithostatic $2225\times9.81\times3200 = \qty{69.9}{MPa}$; hydrostatic
$1030\times9.81\times3200 = \qty{32.3}{MPa}$; \omterm{def:b1:stress-strain-rheology:lithostatic}{effective stress}
$69.9 - 32.3 = \qty{37.5}{MPa}$.

\textbf{6.} $0.49\,e^{-z/3704} = 0.63\,e^{-z/1960}$ gives $z = \ln(0.63/0.49)
/(1/1960 - 1/3704) = \qty{1046}{m}$. Below it the sandstone is the more porous:
shale starts wetter but loses its water over a shorter distance.

\textbf{7.} Sandstone $800 - 0.49\times3704\times(1 - e^{-800/3704}) =
\qty{447.4}{m}$; shale \qty{1306.7}{m}; chalk \qty{546.0}{m}.

\textbf{8.} \qty{2300.2}{m} of grains; water makes $1 - 2300.2/3200 =
\qty{28}{\%}$ of the column.

\textbf{9.} Sandstone $447.4/15 = \qty{29.8}{m}$, shale $1306.7/40 =
\qty{32.7}{m}$, chalk $546.0/20 = \qty{27.3}{m}$ of solid per million years.

\textbf{10.} The shale had the fastest supply, the chalk the slowest --- but the
three differ by only about a fifth, far less than their thicknesses today
suggest.

\textbf{11.} \qty{800}{m}: its top is at the surface today, where its deposition
has just ended, so its thickness at deposition is its thickness today.

\textbf{12.} Solving $s(0, b) = 1306.7$ for the shale by
\cref{met:b1:diagenesis-basins:decompaction} gives $b = \qty{2124}{m}$: the shale
was \qty{2124}{m} thick at \qty{15}{\Ma}.

\textbf{13.} $s(0, b) = 546.0$ for the chalk gives \qty{1071}{m} at
\qty{55}{\Ma}.

\textbf{14.} With its top at \qty{2124}{m}, $s(2124, b) = 546.0$ gives $b =
\qty{2748}{m}$: the chalk was \qty{624}{m} thick.

\textbf{15.} \qty{1071}{m} at \qty{55}{\Ma}; $2124 + 624 = \qty{2748}{m}$ at
\qty{15}{\Ma}.

\textbf{16.} From today's thicknesses: sandstone $800/15 = \qty{53.3}{m}$,
shale $1800/40 = \qty{45.0}{m}$, chalk $600/20 = \qty{30.0}{m}$ per million
years. From the thicknesses at deposition: \qty{53.3}{m}, $2124/40 =
\qty{53.1}{m}$ and $1071/20 = \qty{53.6}{m}$. The second set is the meaningful
one --- the rate at which the sea floor was built up --- and it shows a steady
rate that today's thicknesses hide.

\textbf{17.} $1 - 1800/2124 = \qty{15}{\%}$.

\textbf{18.} The water depth is the same at the start and the end of each
interval, so the \omterm{def:b1:diagenesis-basins:accommodation}{accommodation} created is the growth of the column above the
datum: \qty{1071}{m} from 75 to \qty{55}{\Ma}, $2748 - 1071 = \qty{1677}{m}$ from
55 to \qty{15}{\Ma}, and $3200 - 2748 = \qty{452}{m}$ since \qty{15}{\Ma}.

\textbf{19.} $S = \Delta A - \Delta E$: $1677 - (-60) = \qty{1737}{m}$ from 55 to
\qty{15}{\Ma}; $452 - (-30) = \qty{482}{m}$ since.

\textbf{20.} $1737/40 = \qty{43}{m}$ and $482/15 = \qty{32}{m}$ per million
years.

\textbf{21.} $60/1737 = \qty{3.5}{\%}$ and $30/482 = \qty{6.2}{\%}$. Sea level fell
in both intervals, so \omterm{def:b1:diagenesis-basins:eustasy}{eustasy} hindered the filling: the basement had to subside
that much more to make the same room. \omterm{def:b1:diagenesis-basins:subsidence}{Subsidence} made essentially all of the
\omterm{def:b1:diagenesis-basins:accommodation}{accommodation}.

\textbf{22.} The shale and the chalk were \qty{2748}{m} thick at \qty{15}{\Ma}
and are \qty{2400}{m} today: their \omterm{def:b1:diagenesis-basins:compaction}{compaction} made \qty{348}{m} of room,
\qty{77}{\%} of the \qty{452}{m} filled by the sandstone.

\textbf{23.} Each unit has been thinned by the burial that followed it, by an
amount that depends on its lithology and on how deep it went; the room made in
each interval must be measured with the thicknesses the units had then, which
only \omterm{def:b1:diagenesis-basins:compaction}{decompaction} gives. And even then the \omterm{def:b1:diagenesis-basins:subsidence}{subsidence} so found includes the
sinking caused by the weight of the sediment, which must be removed to reach the
\omterm{def:b1:diagenesis-basins:subsidence}{tectonic subsidence}.

\textbf{24.} \qty{2124}{m}, about \qty{2.1}{km} --- \qty{324}{m} more than the
\qty{1800}{m} the well finds today.
\end{solution}
